% Chương này trình bày bối cảnh, lý do thực hiện, mục tiêu và scope của đề tài.
\chapter{GIỚI THIỆU}

\section{Lý do thực hiện đề tài}
% Chi tiết: Mô tả bối cảnh kinh tế, vấn đề rủi ro tín dụng, và tầm quan trọng của việc dự báo

Trong bối cảnh nền kinh tế số phát triển mạnh mẽ, các tổ chức tài chính và ngân hàng đang chuyển dịch mạnh mẽ từ phương thức thẩm định tín dụng truyền thống sang các hệ thống tự động dựa trên Dữ liệu lớn (Big Data) và Học máy (Machine Learning). Rủi ro tín dụng, cụ thể là khả năng khách hàng không thể hoàn trả nợ đúng hạn, là yếu tố sống còn ảnh hưởng trực tiếp đến tỷ lệ nợ xấu (Non-Performing Loans - NPL) và tính thanh khoản của toàn hệ thống.

Theo số liệu từ Ngân hàng Thế giới, tỷ lệ nợ xấu trung bình toàn cầu dao động ở mức 1-2% tổng tài sản ngân hàng, và ở một số khu vực con số này có thể vượt quá 5%, gây áp lực khổng lồ lên dòng vốn. Do đó, việc xây dựng một hệ thống dự báo vỡ nợ chính xác không chỉ giúp ngân hàng lập rào chắn rủi ro hiệu quả, mà còn cho phép họ mở rộng an toàn tệp khách hàng tiềm năng thông qua việc cá nhân hóa hạn mức tín dụng.

Đồ án này tập trung giải quyết bài toán phân loại nhị phân nhằm dự báo xác suất vỡ nợ của khách hàng thẻ tín dụng. Bằng việc kết hợp các thuật toán học máy tiên tiến thuộc nhóm Ensemble Learning (LightGBM, XGBoost), kỹ thuật cân bằng dữ liệu, và phương pháp tối ưu hóa siêu tham số, nhóm hướng đến việc xây dựng một hệ thống thẩm định tín dụng mạnh mẽ, minh bạch và có tính ứng dụng thực tiễn cao.

\section{Mục tiêu nghiên cứu}

Đề tài được xác định các mục tiêu cụ thể sau:

\begin{enumerate}

    \item \textbf{Phân tích khám phá (EDA) và trích xuất đặc trưng:} Xác định các yếu tố hành vi tác động mạnh nhất đến khả năng vỡ nợ. Ứng dụng Kỹ nghệ đặc trưng (Feature Engineering) để tạo ra các chỉ số tài chính định lượng có sức mạnh dự báo cao như: Tỷ lệ sử dụng tín dụng, Tần suất chậm nợ, Tỷ lệ thanh toán.

    \item \textbf{Xử lý triệt để bài toán mất cân bằng dữ liệu:} Áp dụng kỹ thuật SMOTE (Synthetic Minority Over-sampling Technique) nghiêm ngặt trên tập huấn luyện nhằm giải quyết sự chênh lệch lớn giữa hai lớp khách hàng, đồng thời đảm bảo không xảy ra hiện tượng rò rỉ dữ liệu (data leakage) sang tập kiểm thử.

    \item \textbf{Xây dựng và tối ưu hóa mô hình học máy:} Triển khai đường ống (pipeline) huấn luyện và đánh giá toàn diện 5 mô hình: Logistic Regression (Baseline), LightGBM (Mặc định \& Tối ưu hóa), XGBoost (Mặc định \& Tối ưu hóa). Sử dụng Bayesian Optimization (Optuna với 30 trials) để dò tìm không gian siêu tham số tối ưu.

    \item \textbf{Giải thích mô hình (Explainability):} Ứng dụng giá trị SHAP (SHapley Additive exPlanations) để "mở hộp đen" mô hình, giúp giải thích rõ ràng mức độ đóng góp của từng đặc trưng vào quyết định tín dụng của hệ thống.

\end{enumerate}

\section{Phạm vi và Giới hạn}

\textbf{Phạm vi của đề tài:}
\begin{itemize}
    \item Tập dữ liệu: 30.000 khách hàng thẻ tín dụng từ Đài Loan (Nguồn: UCI Machine Learning Repository).
    \item Phân chia dữ liệu: Chia tập dữ liệu theo tỷ lệ chuẩn 64/16/20 (Train/Validation/Test) có phân tầng (stratify) để giữ nguyên tỷ lệ phân phối nhãn.
    \item Phương pháp: Ensemble Learning (Gradient Boosting) kết hợp cùng Bayesian Optimization.
    \item Metric đánh giá: Sử dụng ROC-AUC làm metric chính, kết hợp cùng F1-Score, Recall, Precision và Confusion Matrix.
\end{itemize}

\textbf{Giới hạn:}
\begin{itemize}
    \item Dữ liệu phản ánh hành vi tài chính tại một khu vực (Đài Loan) và thời điểm tĩnh (2005), có thể chưa phản ánh đầy đủ các biến động kinh tế vĩ mô hiện tại.
    \item Hệ thống tập trung xử lý dữ liệu dạng bảng có cấu trúc (tabular data), chưa bao gồm các dữ liệu phi cấu trúc như lịch sử tín dụng dạng văn bản (text) hay thông tin định danh (images).
    \item Chưa áp dụng cơ chế học trực tuyến (Online Learning) để mô hình tự động cập nhật liên tục với dữ liệu thời gian thực.
\end{itemize}
